%% related -- constrained types and simplification, restricted entailment,
%% feature trees and set constraints, the first-order theory, the automata route
%% to plain entailment, pushdown saturation, algebraic and semantic subtyping,
%% representations of solution spaces, and the proof for plain entailment.

\section{Related Work}\label{sec:rel}

\Cref{tab:rel} summarizes selected results on entailment questions that
involve existential quantifiers or projections.

\begin{table}[tbp]
\caption{Selected results related to existential subtype entailment. ``Open
alphabet'' means that a feature not occurring in the formula is always
available.}
\label{tab:rel}
\small\setlength{\tabcolsep}{3pt}
\begin{tabular}{@{}>{\raggedright\arraybackslash}p{0.35\linewidth}>{\raggedright\arraybackslash}p{0.17\linewidth}>{\raggedright\arraybackslash}p{0.20\linewidth}>{\raggedright\arraybackslash}p{0.22\linewidth}@{}}
\toprule
Problem & Trees & Result & Reference \\
\midrule
Structural subtyping, first-order theory (hence existential entailment) & finite & decidable & \citet{KuncakRinard03} \\
Non-structural subtyping over $\BF$, $f$ binary and covariant, first-order theory & finite, regular, arbitrary & undecidable, with six quantifier blocks & \citet{SuAikenNiehrenPriesnitzTreinen02} \\
Feature-tree ordering constraints, existential entailment, open alphabet & finite, arbitrary & PSPACE-complete & \citet{MullerNiehrenTreinen98,MullerNiehrenTreinen01} \\
Restricted entailment (projection inclusion) in a set-based selector domain & complete infinite & exponential time; PSPACE-hard & \citet{FlanaganFelleisen99} \\
Subtyping of constrained types with $\to$, $\Bot$, $\Top$ & regular & posed as open; sound decidable approximation & \citet{TrifonovSmith96} \\
Subsumption of polar type schemes & extensible algebra & decidable & \citet{Dolan16} \\
Plain \NSSE, one constructor of any arity and variance & finite, regular, arbitrary & decidable & \citet{Zhao26} \\
Existential entailment, one constructor of any arity and variance over $\BFT$, and over $\BF$ with $f$ covariant & finite, regular, arbitrary & decidable; in 2-EXPSPACE; PSPACE-hard for each fixed arity $\ge2$, already with $f$ covariant & this paper \\
\bottomrule
\end{tabular}
\end{table}

\paragraph{Constrained Types and Their Simplification}
Constrained types with recursive constraints were used for type inference for
objects by \citet{EifrigSmithTrifonov95}. \citet{TrifonovSmith96} defined
subtyping between constrained types semantically, over regular trees with
arrows, $\Bot$, and $\Top$; \cref{eq:intro-subsumption} is their
definition, for types without assumptions on program variables, written as an
existential entailment. Their decision procedure is a sound approximation
based on canonical constraint maps and primitive derivations, and
\cref{cor:ap-decidable} decides their relation exactly. \citet{Pottier96}
described an incomplete entailment algorithm, together with the removal of
unreachable constraints and simplifying substitutions. \citet{Pottier01}
developed polarity-based garbage collection, canonization, and minimization,
and proved completeness in the equality setting; completeness is lost for
subtyping. \citet[Section~8.5]{Pottier98thesis} gives a correct simplification
that canonical bounds miss; \cref{ex:ap-pottier} recalls it.
\citet{AikenWimmersPalsberg99} eliminate quantified variables and obtain
optimal representations for unconstrained polymorphic types with subtyping;
their extension to constrained types is sound but not complete.
\citet{FahndrichAiken96} use polarity-based minimization, removal of
intermediate variables, and cycle collapse in an implemented analysis, and they
define the comparison of constrained types by a universal-existential
condition. For constrained types, all these works use sound transformations or
approximations; the classes of \cref{sec:ap-plain} are cheap complete cases, and
\cref{thm:as-main} covers the rest.

\citet{FlanaganFelleisen99} decide a closely related problem. Their restricted
entailment $S_1\entails_E S_2$ is the inclusion $\Sol(S_1)|_E\subseteq\Sol(S_2)|_E$
of projections onto a set $E$ of external variables, which after renaming is
the existential entailment $S_1\entails\exists\bar z\,S_2$ with $\bar z$ the
variables of $S_2$ outside $E$. They decide it completely, in exponential time,
and show PSPACE-hardness. Their domain, however, is the set-based domain
$D=\mathcal P(\mathrm{Const})\times D\times D$ of complete infinite trees with
total selectors and a contravariant domain component, not the ground trees of
non-structural subtyping, so their decision procedure does not apply to our
setting.

\paragraph{Feature Trees and Set Constraints}
For ordering constraints over feature trees, the full first-order theory is
undecidable, while existential entailment is PSPACE-complete over finite and
over possibly infinite trees \citep{MullerNiehrenTreinen98,MullerNiehrenTreinen01}.
Their procedure saturates the constraints, builds polynomial-size automata that
recognize consequences expressed by paths, hides the existential variables
while keeping their states, and decides $\varphi\entails\Exists\psi$ by an
inclusion of the resulting languages. Our three conditions and the pure
derivations through hidden variables play the role of their consequence
languages and hidden states. Their setting differs from ours in one essential
respect: the feature alphabet is open, so a feature that does not occur in the
formula is always available. Their PSPACE-hardness proof, for instance, builds
countermodels with such a fresh feature
\citep[Lemma~5.5]{MullerNiehrenTreinen01}, and \citet{MullerNiehrenPodelski00}
use fresh features to construct countermodels for ordering constraints. A
ranked signature has no fresh direction: a constructor node has exactly its
prescribed children, and $\Bot$ and $\Top$ have none.

The alphabet matters already for feature trees, as the following example over
arbitrary trees shows. Over the features $\{0,1\}$, the constraints
$a(x)\wedge x[0]x\wedge x[1]x\wedge x\le y$ make $x$ the complete infinite tree
labeled $a$ everywhere, so every $y$ above $x$ equals $x$, and the constraints
entail $y\le x$. Once a third feature is available, $y$ may extend $x$ along
it, and the entailment fails. The open-alphabet and the fixed-alphabet
problems thus have different answers on this instance, and a decision
procedure for the first does not decide the second. The same phenomenon
appears in our setting when the arity grows: an entailment for a binary
constructor need not survive bare padding to a larger arity
(\cref{ex:cx-padding}). In set-constraint languages, entailment of atomic set
constraints is PSPACE-complete over infinite signatures
\citep{NiehrenMullerTalbot99}; these constraints have no existential
conclusions, and sets of trees are a different domain from trees under the
subtype order.

\paragraph{The First-Order Theory}
The undecidability proof of \citet{SuAikenNiehrenPriesnitzTreinen02}, recalled
in \cref{sec:intro}, reduces the Post correspondence problem to a sentence whose
matrix uses implications, disjunctions, and negations; for example, it
expresses that the subtypes of a tree form a chain. Our problem is the validity
of $\forall\bar x\,(\varphi\Rightarrow\Exists\psi)$ with conjunctive $\varphi$
and $\psi$, a proper restriction of the $\forall^*\exists^*$ fragment. The same
paper decides the first-order theory over signatures with only unary and
nullary constructors. \citet{Su02} and the extended version of the paper
\citep{SuEtAl06} reduce entailment to emptiness of constrained automata on
tuples of trees, which is undecidable in general. The remark of
\citet[Section~7.2.2]{Su02} quoted in \cref{sec:intro} concerns an adaptation
of that undecidability proof; no formula with a single alternation is
displayed, and undecidability of a fragment with one alternation but an
unrestricted matrix would still not settle the case of conjunctions.
\citet{KuncakRinard03} showed that the first-order theory of structural
subtyping of non-recursive types is decidable, by a term-power construction.
This covers existential entailment for structural subtyping, which compares
only trees of the same shape, over finite trees.

\paragraph{The Automata Route to Plain Entailment}
\citet{NiehrenPriesnitz99} proved PSPACE-completeness of plain non-structural
entailment for the fragment without literals $x=\Bot$ and $x=\Top$, by a
reduction to universality of finite automata.
\citet{NiehrenPriesnitz01,NiehrenPriesnitz03} introduced cap automata,
decided the unary case, and identified the role of word equations in the
general problem. \citet{HengleinRehof98} proved PSPACE-hardness of
non-structural entailment for a binary covariant constructor, which
\cref{thm:cx-hard} extends to every fixed arity at least two, and the theses of
Priesnitz \citep{Priesnitz00,Priesnitz04} develop the automata approach
further. \citet{NiehrenPriesnitzSu05} classify the complexity of
satisfiability over posets and note that their technique does not extend
beyond satisfiability. Our reduction keeps the spirit of this line: the
obstructions are periodicity conditions between suffixes of paths. The new
features are that two words with a common suffix are compared, that a word may
be compared with itself through a self event, and that the separator encoding
returns the two words to the one-word readers of \citet{Zhao26}.

\paragraph{Pushdown Saturation Through Internal Variables}
The type inference of \citet{NoonanLoginovCok16} for machine code uses derived
type variables, which extend a variable by a word of field labels, and observes
that ``the full set of constraints entailed by $C$ has a finite encoding by an
automaton''. It builds a transducer for derivations through uninteresting
variables by a modification of the saturation procedure for prefix rewriting
\citep{Buchi64,Caucal92}, and normalizes derivations to pops followed by
pushes. This is a close precedent for our pure derivations through hidden
variables and their lowest-point decomposition (\cref{lem:sm-lowest-point}).
Their entailment is derivability in a proof system, while ours is semantic
entailment over ground trees. Completeness for the semantic notion is what
needs the leaf-cone characterization and the analysis of two-word holes, which
have no counterpart there. The saturation in our setting is the computation of
the closure $D_\chi$ of the signed system; no further saturation is needed,
because the pushdown summaries are read off from it
(\cref{def:sm-automata,lem:sm-summaries}).

\paragraph{Algebraic and Semantic Subtyping}
Algebraic subtyping \citep{Dolan16,DolanMycroft17} interprets types in an
extensible algebra rather than as ground trees and infers principal types by
biunification. Subsumption between its polar type schemes is decided in the
thesis \citep{Dolan16}; the conference paper does not discuss deciding
subsumption \citep{DolanMycroft17}. Simple-sub \citep{Parreaux20} gives a compact
account of this inference; MLstruct extends it to a Boolean algebra of
structural types and does not formally describe its subsumption check
\citep{ParreauxChau22}; and \citet{ParreauxEtAl24} keep subtyping constraints
in first-class polymorphic types. Because the semantics differ, decidability
there does not decide the ground-tree problem; \cref{app:ap-closed} gives a
correct closed-tree simplification that is not valid in an extensible
algebra. Semantic subtyping interprets types as sets of values.
\citet{CastagnaXu11} decide subtyping between polymorphic regular types in
EXPTIME, and \citet{GesbertGenevesLayaida15} decide semantic subtyping through
a logic. For polymorphic application, \citet{CastagnaNguyenXuAbate15} decide
tallying, but inferring the type substitutions in general uses sound and
complete semi-decision procedures whose termination is left open. These are
questions about types, not about the solutions of constraint systems over
ground trees.

\paragraph{Representations of Solution Spaces}
\citet{Kozen95} represents the solution spaces of set constraints as rational
spaces and relates logical entailment to rational embeddings between these
spaces. Our leaf representation (\cref{thm:lf-representation}) is a
representation theorem of a different kind, in the tradition of rings of sets
\citep{Birkhoff37} and Alexandrov topologies \citep{Alexandroff37}: a tree of
any variance is a pair of disjoint sets of words that are closed under
extension, and the order of trees is inclusion of the first set and reverse
inclusion of the second. This elementary fact becomes useful because every
literal then becomes clauses with a single premise. For a covariant
constructor over $\BF$, solution sets are closed under unions and intersections
(\cref{rem:ch-one-sign}); in general they are closed under intersections of
leaf cones, which in the order of trees is the median with $\Omega$
(\cref{prop:ch-meet}).

\paragraph{Decidability of Plain Entailment}
\Citet{Zhao26} decides plain entailment
$\varphi\entails x\le y$ for one constructor of any arity and variance, by
reducing non-entailment to holes of finite readers and bounding a shortest hole
through its rejected tail, its depth, and a compression argument. It lists
existential entailment among the natural extensions that are not decided by
testing component inequalities independently. We use the results recalled in
\cref{sec:bg} and \cref{app:bg}. Condition \textup{(U)} is a hole problem of
exactly the kind solved there, decided by its three bounds, applied to a
reduced reader whose depth is bounded through a clock on the antecedent
(\cref{sec:un}); the readers of that paper already compare a root with its own
signed copy, the plain counterpart of our self events. For the two-word
holes of condition \textup{(II)}, the rejected tail of a mortal hole is short
by definition, so its depth theorem already bounds the depth of
the separator word (\cref{lem:mo-depth}). Our main technical additions to that
machinery are the pump for holes that are not mortal
(\cref{thm:nm-pump}) and the extension of the bounded-depth search to holes
that must contain an anchored comparison (\cref{thm:mo-anchored}).
